Thermoplastische Kunststoffe finden aufgrund ihrer spezifischen Eigenschaften zunehmend Verwendung als alternative Zahnradwerkstoffe. Die Zahnfußtragfähigkeit von Kunststoffzahnrädern ist ein wesentlicher Faktor für deren Lebensdauer. Beim Einsatz von faserverstärkten Kunststoffen beeinflusst die Ausrichtung der Verstärkungsfasern maßgeblich die mechanischen Festigkeitseigenschaften eines Kunststoffzahnrades. Allerdings existieren bislang keine präzisen Aussagen hinsichtlich der Gültigkeit und der Grenzen der unterschiedlichen Berechnungsverfahren der Zahnfußtragfähigkeit von Kunststoffzahnrädern.

Im Rahmen dieser Semesterarbeit soll die Zahnfußtragfähigkeit kurzfaserverstärkter Kunststoffzahnräder theoretisch und experimentell untersucht werden. Ziel dieser Arbeit ist es, das Potenzial kurzfaserverstärkter Kunststoffzahnräder zur Steigerung der Zahnfußtragfähigkeit zu analysieren. Auf Basis der Ergebnisse sollen die verwendeten Berechnungsmethoden hinsichtlich ihrer Genauigkeit und Anwendbarkeit verglichen und bewertet werden.

Dazu sollen folgende Arbeitspakete bearbeitet werden:
\begin{itemize}
    \item Literaturrecherche zum Stand des Wissens von kurzfaserverstärkten Kunststoffzahnrädern mit Fokus auf deren Tragfähigkeit.
    \item Berechnung der Zahnfußspannung mit der S-Life plastics und FVA Workbench evtl. unter Berücksichtigung optimierter Zahnfußgeometrien (Elliptisch/Bionisch/Bezier).
    \item Berechnung der Zahnfußspannung mit dem FZG Programm RIKOR.
    \item Durchführung von experimentellen Untersuchungen am FZG-kleinverspannungsprüfstand mit Achsabstand $a = 52\,\text{mm}$.
    \item Dokumentation der Prüfräder und Auswertung der Ergebnisse vor und nach den experimentellen Versuchen.
\end{itemize}

%
% Beispiel
% Im Rahmen eines Forschungsvorhabens mit Rolls-Royce plc soll eine neue Triebwerksgeneration auf Basis der Forschungsergebnisse mehrerer Universitäten entwickelt werden. Am Lehrstuhl für Maschinenelemente soll ein Planetengetriebeprüfstand ausgelegt und konstruiert werden, mit welchem sich Dynamik- sowie Wirkungsgraduntersuchungen durchführen lassen. Dabei soll eine Vergleichbarkeit zum größer dimensionierten Originalgetriebe von Rolls-Royce gewährleistet sein. Der Prüfstand wird eine Validierung der Simulationsmodelle, die für das Forschungsprojekt erstellt werden, ermöglichen und neue Erkenntnisse zur Dynamik und zum Verlustverhalten von schnelllaufenden Planetengetrieben liefern.
% Im Rahmen der Masterarbeit sollen die vielfältigen Anforderungen an das Prüfstandskonzept in einem Lastenheft bezüglich der Zielgrößen gegliedert zusammengetragen werden. Auf deren Grundlage soll die Konzepterstellung in Form einer vergleichenden und rechnerischen Gegenüberstellung der möglichen Getriebetopologien für einen am Lehrstuhl einsetzbaren Planetengetriebeprüfstand erfolgen. Das vielversprechendste Konzept soll ausgewählt und konstruktiv umgesetzt werden. Besondere Berücksichtigung sollen dabei die Einhaltung der zugrunde gelegten Kenngrößen sowie die Realisierbarkeit der messtechnischen Erfassung dieser Zielgrößen durch die abzuleitende Prüfstandskonstruktion finden.
% Die Arbeit gliedert sich in folgende Arbeitspakete:
% \begin{itemize}
%   \item Erstellung eines Lastenheftes bezüglich der Anforderungen an den Prüfstand
%   \item Durchführung einer Literaturrecherche zu bestehenden Planetengetriebeprüfständen, dynamischen und wirkungsgradrelevanten Kenngrößen sowie verfügbaren messtechnischen Lösungen 
%   \item Konzepterstellung und -auswahl auf Basis der relevanten Messgrößen, Anforderungen des Lastenheftes und der Gegebenheiten am Lehrstuhl
%   \item Konstruktion und Ausarbeitung des ausgewählten Prüfstandskonzeptes
%   \item Normgerechte Auslegung einer im Rahmen des Konzeptes nutzbaren Verzahnung
%   \item Saubere und vollständige Dokumentation der durchgeführten Arbeiten
% \end{itemize}